\chapter{0. Introduction}
\source{cheeger-gromov-taylor:page-0001,cheeger-gromov-taylor:page-0002,cheeger-gromov-taylor:page-0003,cheeger-gromov-taylor:page-0004}

% Source-faithful transcription of Cheeger--Gromov--Taylor, PDF pages 1--4.

Let $M^n$ be a complete Riemannian manifold. Then the Laplacian $\Delta = -\delta d$
on functions is a nonpositive essentially self adjoint operator when restricted to
functions of compact support. Thus functions $f(\sqrt{-\Delta})$ can be defined by the
spectral theorem for unbounded self adjoint operators, according to the
prescription

\begin{equation}
\tag{0.1}
f(\sqrt{-\Delta}) = \int_0^\infty f(\lambda)\, dE_\lambda,
\end{equation}

where $dE_\lambda$ is the projection valued measure associated with $\sqrt{-\Delta}$.

A natural problem is to study the behavior of the explicit kernel $k_{f(\lambda)}(x_1, x_2)$
representing $f(\sqrt{-\Delta})$, in terms of the behavior of various geometric quantities
on $M^n$. As a particularly important example we have the heat kernel
$E(x_1, x_2, t) = k_{e^{-\lambda^2 t}}$. By use of the local parametrix and the standard elliptic
estimates, one can show that for $t > 0$, $E(x_1, x_2, t)$ is a positive (symmetric)
$C^\infty$ function of $x_1, x_2, t$ which for fixed $t$ and (say) $x_2$, is in the domain of all
positive powers of $\Delta$ as a function of $x_1$; see e.g. [9]. In works of G\"arding [19]
and Donnelly [16], upper estimates for $E(x_1, x_2, t)$ (and its derivatives) were
given under the assumption that $M^n$ has bounded geometry. They showed that
as $x_2 \to \infty$, the behavior of $E(x_1, x_2, t)$ is roughly similar to that of the
Euclidean heat kernel,
\[
\frac{e^{-\rho^2(x_1, x_2)/4}}{(4\pi t)^{n/2}};
\qquad (\rho(x_1, x_2)\text{ denotes distance}).
\]

Recall that

$M^n$ is said to have bounded geometry if the injectivity radius $i(x)$ of the
exponential map at $x$ stays uniformly bounded from below by some constant
$\delta>0$, and if (some or all of) the covariant derivatives $\nabla^i R$ of the
curvature tensor $R$ stay bounded in norm; see \S 3.

Recently, Cheng, Li and Yau [10] achieved a significant advance. They gave
an estimate of the same general character as the above, but under a much
weaker hypothesis. In particular, they showed that one need only assume a
bound on the sectional curvature, $\lvert K_M\rvert < K$, and a lower bound on the
injectivity radius of the exponential map $i(p)$ at some point $p$. Their argument
makes use of three highly nontrivial facts which we now recall.

In what follows we will always let $\lVert \,\cdot\, \rVert_\infty$ denote the sup norm on
$D \subset M^n$, and $\lVert \,\cdot\, \rVert_{L^2(D)}$ the $L^2$ norm on $D$.

\begin{enumerate}
\item \textbf{Moser iteration and the Sobolev constant.} The Sobolev inequality
states the existence of a constant $S_D>0$, such that
\begin{equation}
S_D\left(\int_D \lvert g\rvert^{n/(n-1)}\right)^{(n-1)/n}
\le \int_D \lVert dg\rVert,
\tag{0.2}
\end{equation}
for any $g\in H_0^1(D)$. Let $B_R(x)$ be a metric ball, and let $u(x,t)$ be a solution of
the heat equation $(\Delta-\partial/\partial t)u=0$ on $[0,T]\times B_R(x)$. Then Moser's iteration
argument [30, pp. 115--117] gives the powerful estimate
\begin{equation}
\lvert u(x,T)\rvert \le c_n\left[\frac{1}{T^{(n+2)/2}}+\frac{1}{A^{n+2}}\right]^{1/2}
S_D^{-n/2}\lVert u\rVert_{L^2([0,T]\times B_A(x))}.
\tag{0.3}
\end{equation}
Similarly, if $u$ is harmonic on $[-A/2,A/2]\times B_A(x)$, (($\partial^2/\partial t^2+\Delta)u=0$),
then
\begin{equation}
\lvert u(x,0)\rvert \le c_n A^{-(n+1)/2}S_D^{-n/2}\lVert u\rVert_{L^2([-A/2,A/2]\times B_A(x))}.
\tag{0.3$'$}
\end{equation}

\item \textbf{The Sobolev constant equals the isoperimetric constant.} In [4], [17],
[28], it is shown that for any domain $D$ in a Riemannian manifold $M^n$, the
constant $S_D$ is precisely equal to the optimal constant $\Phi(D)$ in the isoperimetric inequality
\begin{equation}
A(\partial U)\ge \Phi(D)(V(U))^{(n-1)/n},
\tag{0.4}
\end{equation}
which relates the area $A(\partial U)$ and volume $V(U)$ of an arbitrary subdomain
with smooth boundary $U\subset D\subset M^n$.

\item \textbf{Croke's estimate of the isoperimetric constant.} In [13] Croke defines a
constant $\widetilde{\omega}(D)$ as follows. Let $x\in D$, $v\in T_xM^n$, $\lVert v\rVert=1$. Let $\gamma_v$ be the
geodesic with $\gamma'(0)=v$, and let $t_v$ denote the first parameter value $t$, for which
$\gamma_v(t)\in \partial D$. Finally, let $\mu(\,)$ denote angular measure on the unit sphere in
$T_xM^n$, normalized so that the total measure is $1$. Set
\begin{equation}
\widetilde{\omega}(D)=\inf_{x\in D}\mu(\{v\in T_xM^n\mid \lVert v\rVert=1 \text{ and } \gamma_v\lvert_{[0,t_v]} \text{ is minimal}\}).
\tag{0.5}
\end{equation}
\end{enumerate}

Croke proves the striking inequality

\begin{equation}
(0.6)\qquad \Phi(D) \ge 2^{(n-1)/n}\frac{\alpha(n-1)}{\bigl(\alpha(n)\bigr)^{(n-1)/n}}
(\widetilde{\omega}(D))^{(n+1)/n},
\end{equation}

where $\alpha(n)$ is the volume of the unit $n$-sphere; this result is sharp. The proof of
(0.6) depends in an essential way on the fundamental work of Berger [2] and
Kazdan [25].

It is apparent that (0.2), the equality $S_D=\phi(D)$, and (0.6) reduce the
estimation of $E(x_1,x_2,t)$ to bounding the $L^2$-norm $\lVert E(x_1,x_2,t)\rVert_{B_{r_2}(x_2)}$
together with a lower estimate for $\widetilde{\omega}(B_{r_2}(x_2))$. In particular, a considerable
portion of the argument of [10] is devoted to obtaining the required $L^2$-esti-
mate. We enlarge the above circle of ideas by noting that such $L^2$-estimates are
an almost immediate consequence of the finite propagation speed of the wave
kernel $\cos \sqrt{-\Delta_s}\, k_{\cos\lambda s}$. In fact, the statement that $\cos \sqrt{-\Delta_s}$ has unit propa-
gation speed is itself a sharp estimate on $k_{\cos\lambda s}$ which holds universally. To get
estimates on other kernels we need only to synthesize these kernels out of
$k_{\cos\lambda s}$ by means of the Fourier transform; see (1.4). Thus we obtain a
substantial simplification in the derivation of the required $L^2$-estimates, and
extend their scope to a more general class of kernels $k_f$ where $f \in \mathcal{F}_2^{\varphi,A}$; see \S 2
for the definition of $\mathcal{F}_2^{\varphi,A}$.

The pointwise estimates which follow from the previous discussion depend
only on knowing a lower bound for the radius of a ball $B_r(x_0)$ on which we
have a lower bound for $\widetilde{\omega}(B_r(x_0))$. If we define the injectivity angle by

\begin{equation}
(0.7)\qquad \widetilde{\omega}(r,x)=\mu\bigl(\{\gamma'(0)\mid \gamma(0)=x,\ \gamma|_{[0,r]}\text{ is minimal}\}\bigr),
\end{equation}

then clearly

\begin{equation}
(0.8)\qquad \widetilde{\omega}(B_r(x_0)) \ge \inf_{x\in B_r(x_0)} \widetilde{\omega}(2r,x).
\end{equation}

If $r$ is less than or equal to half the injectivity radius $i(x)$, for $x \in B_r(x_0)$ we
have

\begin{equation}
(0.9)\qquad \widetilde{\omega}(B_r(x_0)) = \inf_{x\in B_r(x_0)} \widetilde{\omega}(2r,x)=1.
\end{equation}

In [10] a lower bound for $i(x)$ is established, given $H \le K_M \le K$, a lower
bound on $i(p)$ for some $p$, and the distance $\rho=\rho(p,x)$; they show in
particular that even if $H<0$, $i(x)$ decreases at most exponentially as a
function of $\rho(p,x)$. The injectivity angle $\widetilde{\omega}(2r,x)$ is then estimated using (0.9).
However, a simple comparison argument (see Lemma 4.2) shows that a lower
bound for $\widetilde{\omega}(2r,x)$ follows directly from a lower bound on the Ricci curvature
Ric$_M$ and a lower bound on the volume $V(B_{r_0}(p)) = V_{r_0}(p) > V > 0$; similar
arguments are given in [13], [36] and [10], \S2, Lemma 3, to prove related
statements. From this we obtain (in \S2) an upper estimate on the heat kernel,
and, more generally, on $k_f$ for a large class of functions $f$ as described in \S2,
assuming $\Ric_M \geq (n - 1)H$, $V_r(p) \geq V > 0$.

In fact, it turns out that if a lower bound on the injectivity radius $i(x)$ and a
bound on $|K_M|$ is given, then the appeal to the deep machinery (1)--(3) above
can be avoided. In this case, examination of a parametrix for $\Delta$ leads to an
elliptic estimate, in which the constant can be bounded by a straightforward
comparison argument. The resulting kernel estimates then apply to an essen-
tially wider class of kernels $k_f$, $f \in \mathcal{F}_1^N$, $N \geq n/4$; see \S1. In this connection we
also give an estimate on the decay of $i(x)$, which depends only on assuming
$V(B_{r_0}(p)) \geq V > 0$, rather than $i(p) \geq i_0 > 0$; see Theorem 4.3. Our proof is
based on ideas of [21]; unlike the argument of [10], it does not rely on
Toponogov's theorem.

In dimensions 2 and 3, by means of a special argument we obtain estimates
for the more general functions $f \in \mathcal{F}_1^N$ assuming only $\Ric_M \geq (n - 1)H$,
$V_{r_0}(p) \geq V > 0$. We do not know if these estimates can ultimately be gener-
alized to higher dimensions.

We now briefly summarize the contents of the remaining sections of the
paper.

In \S1 we give the estimates on kernels $k_f(x_1,x_2)$ for $f \in \mathcal{F}_1^N$ assuming a
bound on the sectional curvature, which generalize the results of [10]. In \S2 we
give estimates on kernels $k_f(x_1,x_2)$ for $f \in \mathcal{F}_2^{\varphi,A}$, assuming only a lower
bound on the Ricci tensor. In \S3 we make a more detailed study of various
situations involving bounded geometry. In particular we get bounds on
$k_f(x_1,x_2)$ for certain classes of functions $f$ holomorphic in a strip about the
real axis, satisfying certain ``symbol estimates.'' These bounds yield, in particu-
lar, some $L^p$-operator norm estimates of the sort investigated by Clerc and
Stein [12] and by Stanton and Tomas [31] in the special case of certain classes of
symmetric spaces.

\S4 gives estimates on the volume, injectivity radius, and injectivity angle
used in preceding sections. Here we give special emphasis to the role played by
the relative volume estimates which follow from a lower bound on the Ricci
curvature. As an application we consider manifolds for which the condition
$\Ric_x \geq c/r^2$ holds outside some ball $B_{r_0}(p)$, where $r = \rho(p,x)$. We give sharp
upper and lower bounds on the asymptotic growth of the volume $V_r(p)$ of
$B_r(p)$, in terms of the constant $c$; the qualitative behavior of $V_r(p)$ is strongly
influenced by the precise value of $c$. \S4 can be read independently of the
preceding sections, and some may wish to read it first.
